% !PRINTED 503
\weberpage{503}

Um den Satz vollständig zu begründen, müssen wir noch beweisen, daß die symmetrischen Funktionen der $M_1$, $M_2$, $M_3$ in
dem zuletzt betrachteten Ausnahmefall rationale Funktionen von
$j(\omega)$ sind. Dies erfordert, daß wir untersuchen, was aus (6) wird,
wenn für einen singulären Wert $u$ von $v$ der Nenner $F'(v)$, und,
da $M$ endlich bleibt, auch der Zähler verschwindet. Wir lassen
zunächst wieder $\omega$ variabel und bezeichnen die Wurzeln von
(4) mit

\weq{11}{v_1,\, v_2,\, v_2,\, \ldots,\, v_\nu.}
% [sic] printed repeats "v_2" in the list — presumably an intended index run
% $v_1,\ldots,v_\nu$
Es mögen nun für den betrachteten besonderen Wert von $\omega$

\weq{12}{v_1,\, v_2,\, \ldots,\, v_\lambda}
einander gleich werden, während die übrigen $v_{\lambda+1}, \ldots, v_\nu$ davon
verschieden bleiben. Wir nehmen irgend eine symmetrische Funktion der Größe (12), z. B., für ein unbestimmtes $t$:

\weq{13}{\sigma = (t - v_1),\, (t - v_2),\, \ldots,\, (t - v_\lambda);}
% [sic] printed separates the factors with commas rather than a product sign
wenn wir in $\sigma$ alle Permutationen der Größen (11) ausführen, so
bestimmen wir $r$ Werte

\weq{14}{\sigma_1,\, \sigma_2,\, \ldots,\, \sigma_r,}
wenn
\[
r = \frac{\nu!}{\lambda!\,(\nu - \lambda)!}
\]
die Anzahl der Kombinationen von $\nu$ Größen zu je $\lambda$ (ohne Wiederholung) bedeutet. Über $t$ können wir so verfügen, daß auch für
den singulären Wert $\omega$ \gesperrt{nur eine} der Größen (14) gleich dem
ersten $\sigma$ wird. Die $\sigma$ sind die Wurzeln einer Gleichung
\[
\Phi(x, u) = 0
\]
vom Grade $r$, und für den singulären Wert $x = \sigma$ bleibt $\Phi'(\sigma)$
von Null verschieden.

Betrachten wir nun eine symmetrische Funktion $S$ der (12)
entsprechenden Größe
\[
M_1,\, M_2,\, \ldots,\, M_\lambda,
\]
so nimmt diese durch die $\nu$ Kombinationen der Zahlen $a$, $\partial$, $c$, die
zu den Werten (11) führen, gleichfalls $r$ Werte $S_1, S_2, \ldots, S_r$ an,
und die Summe
\[
\Phi(x, u) \sum \frac{S}{x - \sigma} = \Theta(x, u)
\]
ist eine rationale ganze Funktion von $x$ und $u$. Daraus folgt,
indem man $x$ in $\sigma$ übergehen läßt:
\[
S = \frac{\Theta(\sigma, u)}{\Phi'(u)}.
\]
% [sic] printed denominator $\Phi'(u)$ conflicts with the preceding specialization $x=\sigma$ and should be $\Phi'(\sigma)$
